\documentclass[11pt,a4paper]{article}
\usepackage[utf8]{inputenc}
\usepackage{amsmath,amssymb,amsthm}
\usepackage{listings}
\usepackage{xcolor}
\usepackage[margin=1in]{geometry}
\usepackage{hyperref}

\newtheorem{theorem}{Theorem}

\lstset{
  language=C++,
  basicstyle=\ttfamily\small,
  keywordstyle=\color{blue}\bfseries,
  commentstyle=\color{green!60!black},
  stringstyle=\color{red},
  numbers=left,
  numberstyle=\tiny\color{gray},
  breaklines=true,
  frame=single,
  tabsize=4,
  showstringspaces=false
}

\title{IOI 2003 -- Boundary (Convex Hull)}
\author{Solution}
\date{}

\begin{document}
\maketitle

\section{Problem Statement Summary}
Given $N$ points ($N \le 10^5$) in the plane, compute the convex hull.
Output the hull vertices in order, together with the perimeter, area,
and the number of input points lying on the hull boundary.

\section{Solution: Andrew's Monotone Chain}

\subsection{Algorithm}
\begin{enumerate}
  \item Sort points lexicographically by $(x, y)$.
  \item \textbf{Lower hull:} Scan left to right, maintaining a stack.
        Pop the top while the last two stack points and the new point
        make a clockwise (or collinear) turn.
  \item \textbf{Upper hull:} Scan right to left with the same logic.
  \item Concatenate the two hulls (removing duplicated endpoints).
\end{enumerate}

\subsection{Cross Product Test}
For points $A, B, C$, the cross product
$(B - A) \times (C - A) = (B_x - A_x)(C_y - A_y) - (B_y - A_y)(C_x - A_x)$
is positive for a counter-clockwise turn, zero for collinear, and
negative for clockwise.

For a \emph{strict} convex hull (no collinear points on edges), pop
when the cross product is $\le 0$. To include collinear boundary
points, pop only when strictly $< 0$.

\section{C++ Implementation}

\begin{lstlisting}
#include <bits/stdc++.h>
using namespace std;
typedef long long ll;

struct Point {
    ll x, y;
    Point(ll x = 0, ll y = 0) : x(x), y(y) {}
    Point operator-(const Point& o) const { return {x - o.x, y - o.y}; }
    ll cross(const Point& o) const { return x * o.y - y * o.x; }
    bool operator<(const Point& o) const {
        return x < o.x || (x == o.x && y < o.y);
    }
    bool operator==(const Point& o) const {
        return x == o.x && y == o.y;
    }
};

// Returns convex hull in CCW order.
// If includeBoundary is true, collinear points on edges are kept.
vector<Point> convexHull(vector<Point> pts, bool includeBoundary = false) {
    int n = pts.size();
    if (n < 2) return pts;
    sort(pts.begin(), pts.end());
    pts.erase(unique(pts.begin(), pts.end()), pts.end());
    n = pts.size();
    if (n < 2) return pts;

    auto bad = [&](ll v) -> bool {
        return includeBoundary ? (v < 0) : (v <= 0);
    };

    vector<Point> hull;

    // Lower hull
    for (int i = 0; i < n; i++) {
        while (hull.size() >= 2) {
            Point a = hull[hull.size() - 2], b = hull.back();
            if (bad((b - a).cross(pts[i] - a)))
                hull.pop_back();
            else break;
        }
        hull.push_back(pts[i]);
    }

    // Upper hull
    int lower_size = hull.size();
    for (int i = n - 2; i >= 0; i--) {
        while ((int)hull.size() > lower_size) {
            Point a = hull[hull.size() - 2], b = hull.back();
            if (bad((b - a).cross(pts[i] - a)))
                hull.pop_back();
            else break;
        }
        hull.push_back(pts[i]);
    }
    hull.pop_back(); // remove duplicate of first point
    return hull;
}

double perimeter(const vector<Point>& hull) {
    double peri = 0;
    int n = hull.size();
    for (int i = 0; i < n; i++) {
        Point d = hull[(i + 1) % n] - hull[i];
        peri += sqrt((double)(d.x * d.x + d.y * d.y));
    }
    return peri;
}

double area(const vector<Point>& hull) {
    ll a = 0;
    int n = hull.size();
    for (int i = 0; i < n; i++)
        a += hull[i].cross(hull[(i + 1) % n]);
    return abs(a) / 2.0;
}

// Count input points on the hull boundary.
int countOnBoundary(const vector<Point>& hull,
                    const vector<Point>& allPts) {
    int cnt = 0;
    int h = hull.size();
    for (auto& p : allPts) {
        for (int i = 0; i < h; i++) {
            Point a = hull[i], b = hull[(i + 1) % h];
            Point ab = b - a, ap = p - a;
            if (ab.cross(ap) == 0 &&
                min(a.x, b.x) <= p.x && p.x <= max(a.x, b.x) &&
                min(a.y, b.y) <= p.y && p.y <= max(a.y, b.y)) {
                cnt++;
                break;
            }
        }
    }
    return cnt;
}

int main() {
    ios::sync_with_stdio(false);
    cin.tie(nullptr);

    int N;
    cin >> N;
    vector<Point> pts(N);
    for (int i = 0; i < N; i++)
        cin >> pts[i].x >> pts[i].y;

    vector<Point> hull = convexHull(pts, false);

    cout << hull.size() << "\n";
    for (auto& p : hull)
        cout << p.x << " " << p.y << "\n";

    cout << fixed << setprecision(2) << perimeter(hull) << "\n";
    cout << fixed << setprecision(1) << area(hull) << "\n";
    cout << countOnBoundary(hull, pts) << "\n";

    return 0;
}
\end{lstlisting}

\section{Complexity Analysis}
\begin{itemize}
  \item \textbf{Convex hull:} $O(N \log N)$ (sorting dominates; the
        scan itself is $O(N)$ since each point is pushed and popped at
        most once).
  \item \textbf{Boundary count:} $O(N \cdot H)$ where $H$ is the hull
        size. For large inputs, this can be improved to $O(N \log H)$
        by binary-searching the hull edge for each point.
  \item \textbf{Space:} $O(N)$.
\end{itemize}

\end{document}
